\section{Correctness Verification}


\subsection{Round-Trip Property}

The most important correctness test: $\text{INTT}(\text{NTT}(a)) = a$ for all polynomials $a$.

\begin{lstlisting}[language=Go,caption={NTT round-trip fuzz test}]
func FuzzNTTRoundTrip(f *testing.F) {
    f.Add(make([]byte, 512))
    f.Fuzz(func(t *testing.T, data []byte) {
        if len(data) < 512 {
            t.Skip()
        }
        var a, original [256]int32
        for i := range a {
            a[i] = int32(binary.LittleEndian.Uint16(
                data[2*i:])) % Q
            original[i] = a[i]
        }
        NTT(&a)
        INTT(&a)
        for i := range a {
            got := reduce(a[i])
            want := reduce(original[i])
            if got != want {
                t.Fatalf("index %d: got %d, want %d", i, got, want)
            }
        }
    })
}
\end{lstlisting}

\subsection{Homomorphism Property}

$\text{NTT}(a \cdot b) = \text{NTT}(a) \otimes \text{NTT}(b)$ where $\otimes$ is
pointwise multiplication. This verifies that the NTT correctly enables fast polynomial
multiplication.

\subsection{Cross-Implementation Testing}

We test our Go+AVX2 implementation against the reference C implementation (pqcrystals)
and a pure Python implementation. All three must produce identical outputs for the
NIST known-answer test vectors.

%% ============================================================
